\section{Methodology}
\label{sec:methodology}

Our evaluation methodology measures attacker effectiveness, shield safety, defender robustness, and adaptation gains in a closed-loop CPS setting. Figure~\ref{fig:pipeline} illustrates the experiment pipeline from action generation through execution and detection.

\begin{figure}[t]
    \centering
    \begin{tikzpicture}[
        >=Stealth,
        node distance=0.35cm,
        every node/.style={font=\scriptsize},
        step/.style={draw, rounded corners=2pt, minimum width=2.8cm, minimum height=0.5cm, align=center, thick},
        decision/.style={draw, diamond, aspect=2.2, minimum width=1.5cm, inner sep=1pt, align=center, thick},
        arr/.style={->, thick, >=Stealth},
    ]

    %% Pipeline steps
    \node[step, fill=gray!10] (init) {1. Initialise PLC \&\\Factory~I/O scene};
    \node[step, fill=gray!10, below=of init] (warmup) {2. Warm-up period\\(baseline telemetry)};
    \node[step, fill=red!10, below=of warmup] (propose) {3. Attacker proposes\\structured action(s)};
    \node[step, fill=yellow!15, below=of propose] (validate) {4. JSON schema\\validation \& normalise};
    \node[decision, fill=yellow!20, below=0.45cm of validate] (shield) {Shield\\check?};
    \node[step, fill=green!10, below left=0.5cm and -0.3cm of shield] (exec) {5. Execute on\\live PLC};
    \node[step, fill=red!5, below right=0.5cm and -0.3cm of shield] (reject) {Reject \&\\log reason};
    \node[step, fill=blue!10, below=0.8cm of shield] (detect) {6. Defenders\\observe \& score};
    \node[decision, fill=blue!15, below=0.45cm of detect] (detected) {Attack\\detected?};
    \node[step, fill=green!10, below left=0.45cm and -0.3cm of detected] (alert) {Emit\\alert};
    \node[step, fill=orange!10, below right=0.45cm and -0.3cm of detected] (hardcase) {Hard case\\$\to$ bank};
    \node[step, fill=gray!10, below=0.8cm of detected] (log) {7. Record artifacts\\(Parquet, JSON)};

    %% Arrows
    \draw[arr] (init) -- (warmup);
    \draw[arr] (warmup) -- (propose);
    \draw[arr] (propose) -- (validate);
    \draw[arr] (validate) -- (shield);
    \draw[arr] (shield) -- node[left, font=\tiny]{\checkmark} (exec);
    \draw[arr] (shield) -- node[right, font=\tiny]{$\times$} (reject);
    \draw[arr] (exec) |- (detect);
    \draw[arr] (detect) -- (detected);
    \draw[arr] (detected) -- node[left, font=\tiny]{yes} (alert);
    \draw[arr] (detected) -- node[right, font=\tiny]{no} (hardcase);
    \draw[arr] (alert) |- (log);
    \draw[arr] (hardcase) |- (log);

    \end{tikzpicture}
    \caption{Batch-mode experiment pipeline. Each proposed action passes through schema validation and the safety shield before PLC execution. Undetected attacks are stored as hard cases for adaptation.}
    \label{fig:pipeline}
\end{figure}

\subsection{Experiment Workflow}
Each experiment proceeds as follows:
\begin{enumerate}[leftmargin=1.2em]
    \item Initialise the process in a known safe state.
    \item Start the PLC--Factory~I/O control loop and the logging service.
    \item Execute a benign warm-up period to establish baseline telemetry.
    \item Invoke the attacker to propose one or more structured actions.
    \item Pass each action through the safety shield.
    \item Execute approved actions and continue polling.
    \item Run defenders online over the collected snapshots.
    \item Record attack outcome, physical impact, detector behaviour, and shield events.
    \item Store the run as a replayable artifact.
\end{enumerate}

In agent mode, steps 4--7 are replaced by continuous concurrent operation: the attacker and defender agents independently observe plant state and propose actions, all mediated by the coordinator.

\subsection{Attacker Variants}
We compare four attacker classes:
\begin{itemize}[leftmargin=1.2em]
    \item \textbf{Scripted attacker:} deterministic templates with targeted process-aware actions.
    \item \textbf{Random attacker:} bounded random perturbations over the configured attack surface.
    \item \textbf{LLM attacker (batch):} process-aware planning over structured scene and state context, executed through the batch orchestration pipeline.
    \item \textbf{LLM attacker agent:} an autonomous LLM agent that observes plant state in real time, maintains action history with feedback, and proposes attacks concurrently with the defender agent.
\end{itemize}

\subsection{Defender Variants}
We evaluate four defender configurations:
\begin{itemize}[leftmargin=1.2em]
    \item \textbf{Threshold detector:} configurable per-tag thresholds with severity levels.
    \item \textbf{Invariant detector:} rule-based state-consistency checks derived from scene-specific safety properties.
    \item \textbf{Sequence detector:} a sliding-window ML model for temporal anomaly detection.
    \item \textbf{LLM defender agent:} an autonomous LLM agent (agent mode only) that analyses process telemetry, diagnoses anomalies, and proposes corrective actions.
\end{itemize}

\subsection{Metrics}
We measure:
\begin{itemize}[leftmargin=1.2em]
    \item \textbf{Attack validity:} fraction of attacker outputs that compile and execute successfully.
    \item \textbf{Attack success:} fraction of attacks that produce measurable process deviation.
    \item \textbf{Physical impact:} process deviation measured as the difference between attack-period and baseline tag values.
    \item \textbf{Shield effectiveness:} fraction of proposals approved versus rejected, and unsafe action block rate.
    \item \textbf{Detection performance:} precision, recall, F1, false-positive count, and detection latency.
    \item \textbf{Adaptation gain:} improvement in detection robustness across iterative rounds.
\end{itemize}

For agent mode experiments, we additionally report:
\begin{itemize}[leftmargin=1.2em]
    \item \textbf{Agent actions per round:} number of attack and corrective actions proposed per run.
    \item \textbf{Corrective success rate:} fraction of defender corrective actions that restore process state.
    \item \textbf{Adversarial dynamics:} temporal interplay between attacker perturbations and defender corrections.
\end{itemize}

\subsection{Scenes and Scale}

The evaluation covers three Factory~I/O scenes of increasing complexity: From A to B (4 tags, discrete transport), Level Control (15 tags, continuous process), and Sorting by Height (30 tags, mixed discrete/continuous). For batch mode, each attacker variant is evaluated per scene. For agent mode, LLM attacker and LLM defender agents are evaluated on each scene. Each configuration runs 120--200 steps at a 500\,ms polling interval against the live PLC.

\subsection{Closed-Loop Adaptation}
After each evaluation round, traces corresponding to missed or weakly detected attacks are added to a hard-case bank. These traces are used to refine learned defenders or detector policies in the next round. We measure adaptation gain as the improvement in detection F1 between the initial round and subsequent rounds trained on accumulated hard cases.

\subsection{Ground Truth}
\label{sec:methodology:gt}
Ground-truth labelling is derived from the experiment protocol itself. Each polling step is labelled as \emph{attack-active} or \emph{normal} based on whether a shield-approved attack action was active during that step. This label is recorded in the trace metadata at runtime, not post hoc. Baseline normal-operation traces (no attacker active) are collected for each scene to provide unlabelled ground truth for anomaly-detection calibration.